% =====================================================================
% TECT canonical LaTeX preamble (revtex4-2 / single-column / theorem-safe)
% =====================================================================
% Maintainer: Jusang Lee (jtkor@outlook.com)
% Binding from: 2026-05-06
% Purpose: a single source-of-truth preamble for every TECT paper
% (Paper-00 .. Paper-10 + Auxiliary notes). Designed to (a) eliminate
% font-corruption on pdflatex (T1 fontenc + lmodern), (b) make
% theorem-heavy mathematical-physics content render cleanly in a
% single-column layout (PRD class option, NOT PRL two-column), (c)
% provide consistent theorem numbering across all TECT papers via
% amsthm with a shared counter set, (d) make hyperref + cleveref
% cross-references resolve cleanly with the standard two-pass
% pdflatex compile.
%
% Usage in each Paper-NN.tex:
%   % =====================================================================
% TECT canonical LaTeX preamble (revtex4-2 / single-column / theorem-safe)
% =====================================================================
% Maintainer: Jusang Lee (jtkor@outlook.com)
% Binding from: 2026-05-06
% Purpose: a single source-of-truth preamble for every TECT paper
% (Paper-00 .. Paper-10 + Auxiliary notes). Designed to (a) eliminate
% font-corruption on pdflatex (T1 fontenc + lmodern), (b) make
% theorem-heavy mathematical-physics content render cleanly in a
% single-column layout (PRD class option, NOT PRL two-column), (c)
% provide consistent theorem numbering across all TECT papers via
% amsthm with a shared counter set, (d) make hyperref + cleveref
% cross-references resolve cleanly with the standard two-pass
% pdflatex compile.
%
% Usage in each Paper-NN.tex:
%   % =====================================================================
% TECT canonical LaTeX preamble (revtex4-2 / single-column / theorem-safe)
% =====================================================================
% Maintainer: Jusang Lee (jtkor@outlook.com)
% Binding from: 2026-05-06
% Purpose: a single source-of-truth preamble for every TECT paper
% (Paper-00 .. Paper-10 + Auxiliary notes). Designed to (a) eliminate
% font-corruption on pdflatex (T1 fontenc + lmodern), (b) make
% theorem-heavy mathematical-physics content render cleanly in a
% single-column layout (PRD class option, NOT PRL two-column), (c)
% provide consistent theorem numbering across all TECT papers via
% amsthm with a shared counter set, (d) make hyperref + cleveref
% cross-references resolve cleanly with the standard two-pass
% pdflatex compile.
%
% Usage in each Paper-NN.tex:
%   % =====================================================================
% TECT canonical LaTeX preamble (revtex4-2 / single-column / theorem-safe)
% =====================================================================
% Maintainer: Jusang Lee (jtkor@outlook.com)
% Binding from: 2026-05-06
% Purpose: a single source-of-truth preamble for every TECT paper
% (Paper-00 .. Paper-10 + Auxiliary notes). Designed to (a) eliminate
% font-corruption on pdflatex (T1 fontenc + lmodern), (b) make
% theorem-heavy mathematical-physics content render cleanly in a
% single-column layout (PRD class option, NOT PRL two-column), (c)
% provide consistent theorem numbering across all TECT papers via
% amsthm with a shared counter set, (d) make hyperref + cleveref
% cross-references resolve cleanly with the standard two-pass
% pdflatex compile.
%
% Usage in each Paper-NN.tex:
%   \input{../_shared/tect-preamble.tex}
%   \begin{document}
%   ...
%   \end{document}
%
% Build (every paper): `pdflatex Paper-NN && pdflatex Paper-NN`
% (the second pass resolves \ref{} / \cref{} forward references).
% A bash/PowerShell wrapper that automates this is at
% Docs/papers/papers/_shared/build-paper.sh / build-paper.ps1.
% =====================================================================

% ---- Document class --------------------------------------------------
% PRD class option of revtex4-2 with [reprint,onecolumn] gives the same
% APS heading / by-line / abstract style as PRL but in a wider single
% column that handles long display equations and theorem environments
% without column-break artefacts. nofootinbib keeps citations out of
% footnotes. The 11pt option restores standard font metrics; with T1
% fontenc + lmodern this gives non-bitmap glyphs.
\documentclass[%
  aps,prd,reprint,onecolumn,nofootinbib,11pt,%
  amsmath,amssymb,floatfix,longbibliography%
]{revtex4-2}

% ---- Encoding & fonts ------------------------------------------------
\usepackage[T1]{fontenc}        % proper 8-bit font encoding (no font corruption)
\usepackage[utf8]{inputenc}     % allow UTF-8 source
\usepackage{lmodern}            % scalable Latin Modern (replaces bitmap CM)
\usepackage{textcomp}           % extra text symbols
\usepackage{microtype}          % subtle kerning improvements

% ---- UTF-8 → LaTeX character mappings --------------------------------
% Maps the Unicode characters that occur in TECT body text (legacy from
% the chat-content auto-archival rule, CLAUDE.md §4) to LaTeX-safe
% commands. Without these declarations, T1+inputenc rejects the chars
% with "Unicode character X not set up for use with LaTeX".
% Adding declarations centrally (here) is preferable to per-file edits.
\DeclareUnicodeCharacter{00A7}{\S}              % §  (section sign)
\DeclareUnicodeCharacter{00B2}{\ensuremath{^{2}}}        % ²
\DeclareUnicodeCharacter{00B3}{\ensuremath{^{3}}}        % ³
\DeclareUnicodeCharacter{00B5}{\ensuremath{\mu}}         % µ (micro)
\DeclareUnicodeCharacter{00B7}{\ensuremath{\cdot}}       % · (middle dot)
\DeclareUnicodeCharacter{00D7}{\ensuremath{\times}}      % ×
\DeclareUnicodeCharacter{00E4}{\"a}             % ä
\DeclareUnicodeCharacter{00E9}{\'e}             % é
\DeclareUnicodeCharacter{00F6}{\"o}             % ö
\DeclareUnicodeCharacter{00FC}{\"u}             % ü
\DeclareUnicodeCharacter{010C}{\v{C}}           % Č
\DeclareUnicodeCharacter{010D}{\v{c}}           % č
\DeclareUnicodeCharacter{0394}{\ensuremath{\Delta}}      % Δ
\DeclareUnicodeCharacter{03A3}{\ensuremath{\Sigma}}      % Σ
\DeclareUnicodeCharacter{03A9}{\ensuremath{\Omega}}      % Ω
\DeclareUnicodeCharacter{03B1}{\ensuremath{\alpha}}      % α
\DeclareUnicodeCharacter{03B2}{\ensuremath{\beta}}       % β
\DeclareUnicodeCharacter{03B3}{\ensuremath{\gamma}}      % γ
\DeclareUnicodeCharacter{03B4}{\ensuremath{\delta}}      % δ
\DeclareUnicodeCharacter{03B5}{\ensuremath{\varepsilon}} % ε
\DeclareUnicodeCharacter{03BA}{\ensuremath{\kappa}}      % κ
\DeclareUnicodeCharacter{03BB}{\ensuremath{\lambda}}     % λ
\DeclareUnicodeCharacter{03BC}{\ensuremath{\mu}}         % μ
\DeclareUnicodeCharacter{03BD}{\ensuremath{\nu}}         % ν
\DeclareUnicodeCharacter{03C0}{\ensuremath{\pi}}         % π
\DeclareUnicodeCharacter{03C1}{\ensuremath{\rho}}        % ρ
\DeclareUnicodeCharacter{03C3}{\ensuremath{\sigma}}      % σ
\DeclareUnicodeCharacter{03C4}{\ensuremath{\tau}}        % τ
\DeclareUnicodeCharacter{03C6}{\ensuremath{\varphi}}     % φ
\DeclareUnicodeCharacter{03C7}{\ensuremath{\chi}}        % χ
\DeclareUnicodeCharacter{03C8}{\ensuremath{\psi}}        % ψ
\DeclareUnicodeCharacter{03C9}{\ensuremath{\omega}}      % ω
\DeclareUnicodeCharacter{2013}{--}              % – (en dash)
\DeclareUnicodeCharacter{2014}{---}             % — (em dash)
\DeclareUnicodeCharacter{2018}{`}               % ‘
\DeclareUnicodeCharacter{2019}{'}               % ’
\DeclareUnicodeCharacter{201C}{``}              % “
\DeclareUnicodeCharacter{201D}{''}              % ”
\DeclareUnicodeCharacter{2070}{\ensuremath{^{0}}}        % ⁰
\DeclareUnicodeCharacter{2071}{\ensuremath{^{i}}}        % ⁱ
\DeclareUnicodeCharacter{2122}{\textsuperscript{TM}}     % ™
\DeclareUnicodeCharacter{2126}{\ensuremath{\Omega}}      % Ω (ohm sign)
\DeclareUnicodeCharacter{210F}{\ensuremath{\hbar}}       % ℏ
\DeclareUnicodeCharacter{2115}{\ensuremath{\mathbb{N}}}  % ℕ
\DeclareUnicodeCharacter{2102}{\ensuremath{\mathbb{C}}}  % ℂ
\DeclareUnicodeCharacter{211D}{\ensuremath{\mathbb{R}}}  % ℝ
\DeclareUnicodeCharacter{2124}{\ensuremath{\mathbb{Z}}}  % ℤ
\DeclareUnicodeCharacter{211A}{\ensuremath{\mathbb{Q}}}  % ℚ
\DeclareUnicodeCharacter{2192}{\ensuremath{\to}}         % →
\DeclareUnicodeCharacter{21D2}{\ensuremath{\Rightarrow}} % ⇒
\DeclareUnicodeCharacter{2200}{\ensuremath{\forall}}     % ∀
\DeclareUnicodeCharacter{2203}{\ensuremath{\exists}}     % ∃
\DeclareUnicodeCharacter{2208}{\ensuremath{\in}}         % ∈
\DeclareUnicodeCharacter{2209}{\ensuremath{\notin}}      % ∉
\DeclareUnicodeCharacter{2227}{\ensuremath{\wedge}}      % ∧
\DeclareUnicodeCharacter{2228}{\ensuremath{\vee}}        % ∨
\DeclareUnicodeCharacter{2229}{\ensuremath{\cap}}        % ∩
\DeclareUnicodeCharacter{222A}{\ensuremath{\cup}}        % ∪
\DeclareUnicodeCharacter{2245}{\ensuremath{\cong}}       % ≅
\DeclareUnicodeCharacter{2248}{\ensuremath{\approx}}     % ≈
\DeclareUnicodeCharacter{2260}{\ensuremath{\neq}}        % ≠
\DeclareUnicodeCharacter{2264}{\ensuremath{\leq}}        % ≤
\DeclareUnicodeCharacter{2265}{\ensuremath{\geq}}        % ≥
\DeclareUnicodeCharacter{22C5}{\ensuremath{\cdot}}       % ⋅
\DeclareUnicodeCharacter{2295}{\ensuremath{\oplus}}      % ⊕
\DeclareUnicodeCharacter{2297}{\ensuremath{\otimes}}     % ⊗

% ---- Mathematics & symbols ------------------------------------------
\usepackage{amsmath,amssymb,bm,mathrsfs,mathtools}
\usepackage{amsthm}             % theorem environments (load BEFORE hyperref)

% ---- Theorem environments (shared counter set, TECT canonical) -------
% All theorems / lemmas / propositions / corollaries / remarks share a
% single counter so cross-references read consistently across a paper.
% Definitions get a separate counter (definitions are numbered in their
% own sequence by editorial convention).
\newtheorem{theorem}{Theorem}
\newtheorem{lemma}[theorem]{Lemma}
\newtheorem{proposition}[theorem]{Proposition}
\newtheorem{corollary}[theorem]{Corollary}

\theoremstyle{remark}
\newtheorem{remark}[theorem]{Remark}
\newtheorem{example}[theorem]{Example}

\theoremstyle{definition}
\newtheorem{definition}[theorem]{Definition}
\newtheorem{conjecture}[theorem]{Conjecture}
\newtheorem{hypothesis}[theorem]{Hypothesis}

% ---- Tables / floats / graphics --------------------------------------
\usepackage{booktabs}           % \toprule / \midrule / \bottomrule
\usepackage{array}
\usepackage{graphicx}
\usepackage{enumitem}           % \begin{itemize}[leftmargin=...] options
\usepackage{hyphenat}           % \hyp{} for breakable hyphens in \texttt
\usepackage{seqsplit}           % \seqsplit{} for long unbreakable strings
\sloppy                         % allow stretchier inter-word spacing to
                                % reduce overfull \hbox warnings on long URLs
                                % and long \texttt tokens (paragraph-level).
\emergencystretch=3em           % last-resort hbox stretch budget

% ---- Cross-references (hyperref last among the standard packages,
%      cleveref AFTER hyperref) ----------------------------------------
\usepackage[%
  colorlinks=true,%
  linkcolor=blue,%
  citecolor=blue,%
  urlcolor=blue,%
  breaklinks=true,%
  bookmarks=true,%
  bookmarksnumbered=true%
]{hyperref}
\usepackage[capitalise,nameinlink]{cleveref}

% ---- TECT-canonical author block macros -----------------------------
% (defined here so every paper has the same author / affiliation style)
\newcommand{\tectauthor}{%
  \author{Jusang Lee}%
  \email{jtkor@outlook.com}%
  \affiliation{Independent researcher, Republic of Korea}%
}

% ---- TECT-canonical archive footer ----------------------------------
\newcommand{\tectarchivefooter}{%
  \par\medskip
  \noindent\small\textit{Canonical archive}:
  \url{https://github.com/TECT-OpenPhysics/TECT}.\par%
}

% ---- TECT TODO / AUDIT-FLAG / HONEST-SCOPE inline marks --------------
% Used in drafts to highlight pending audit items without disturbing
% the body text style. Suppress via \tectsuppressmarks (define empty).
\providecommand{\tectauditflag}[1]{\textcolor{red}{\textbf{[AUDIT-FLAG: #1]}}}
\providecommand{\tecttodo}[1]{\textcolor{orange}{\textbf{[TODO: #1]}}}
\providecommand{\tecthonestscope}[1]{\textit{Honest scope: #1.}}

% =====================================================================
% End of TECT canonical preamble
% =====================================================================

%   \begin{document}
%   ...
%   \end{document}
%
% Build (every paper): `pdflatex Paper-NN && pdflatex Paper-NN`
% (the second pass resolves \ref{} / \cref{} forward references).
% A bash/PowerShell wrapper that automates this is at
% Docs/papers/papers/_shared/build-paper.sh / build-paper.ps1.
% =====================================================================

% ---- Document class --------------------------------------------------
% PRD class option of revtex4-2 with [reprint,onecolumn] gives the same
% APS heading / by-line / abstract style as PRL but in a wider single
% column that handles long display equations and theorem environments
% without column-break artefacts. nofootinbib keeps citations out of
% footnotes. The 11pt option restores standard font metrics; with T1
% fontenc + lmodern this gives non-bitmap glyphs.
\documentclass[%
  aps,prd,reprint,onecolumn,nofootinbib,11pt,%
  amsmath,amssymb,floatfix,longbibliography%
]{revtex4-2}

% ---- Encoding & fonts ------------------------------------------------
\usepackage[T1]{fontenc}        % proper 8-bit font encoding (no font corruption)
\usepackage[utf8]{inputenc}     % allow UTF-8 source
\usepackage{lmodern}            % scalable Latin Modern (replaces bitmap CM)
\usepackage{textcomp}           % extra text symbols
\usepackage{microtype}          % subtle kerning improvements

% ---- UTF-8 → LaTeX character mappings --------------------------------
% Maps the Unicode characters that occur in TECT body text (legacy from
% the chat-content auto-archival rule, CLAUDE.md §4) to LaTeX-safe
% commands. Without these declarations, T1+inputenc rejects the chars
% with "Unicode character X not set up for use with LaTeX".
% Adding declarations centrally (here) is preferable to per-file edits.
\DeclareUnicodeCharacter{00A7}{\S}              % §  (section sign)
\DeclareUnicodeCharacter{00B2}{\ensuremath{^{2}}}        % ²
\DeclareUnicodeCharacter{00B3}{\ensuremath{^{3}}}        % ³
\DeclareUnicodeCharacter{00B5}{\ensuremath{\mu}}         % µ (micro)
\DeclareUnicodeCharacter{00B7}{\ensuremath{\cdot}}       % · (middle dot)
\DeclareUnicodeCharacter{00D7}{\ensuremath{\times}}      % ×
\DeclareUnicodeCharacter{00E4}{\"a}             % ä
\DeclareUnicodeCharacter{00E9}{\'e}             % é
\DeclareUnicodeCharacter{00F6}{\"o}             % ö
\DeclareUnicodeCharacter{00FC}{\"u}             % ü
\DeclareUnicodeCharacter{010C}{\v{C}}           % Č
\DeclareUnicodeCharacter{010D}{\v{c}}           % č
\DeclareUnicodeCharacter{0394}{\ensuremath{\Delta}}      % Δ
\DeclareUnicodeCharacter{03A3}{\ensuremath{\Sigma}}      % Σ
\DeclareUnicodeCharacter{03A9}{\ensuremath{\Omega}}      % Ω
\DeclareUnicodeCharacter{03B1}{\ensuremath{\alpha}}      % α
\DeclareUnicodeCharacter{03B2}{\ensuremath{\beta}}       % β
\DeclareUnicodeCharacter{03B3}{\ensuremath{\gamma}}      % γ
\DeclareUnicodeCharacter{03B4}{\ensuremath{\delta}}      % δ
\DeclareUnicodeCharacter{03B5}{\ensuremath{\varepsilon}} % ε
\DeclareUnicodeCharacter{03BA}{\ensuremath{\kappa}}      % κ
\DeclareUnicodeCharacter{03BB}{\ensuremath{\lambda}}     % λ
\DeclareUnicodeCharacter{03BC}{\ensuremath{\mu}}         % μ
\DeclareUnicodeCharacter{03BD}{\ensuremath{\nu}}         % ν
\DeclareUnicodeCharacter{03C0}{\ensuremath{\pi}}         % π
\DeclareUnicodeCharacter{03C1}{\ensuremath{\rho}}        % ρ
\DeclareUnicodeCharacter{03C3}{\ensuremath{\sigma}}      % σ
\DeclareUnicodeCharacter{03C4}{\ensuremath{\tau}}        % τ
\DeclareUnicodeCharacter{03C6}{\ensuremath{\varphi}}     % φ
\DeclareUnicodeCharacter{03C7}{\ensuremath{\chi}}        % χ
\DeclareUnicodeCharacter{03C8}{\ensuremath{\psi}}        % ψ
\DeclareUnicodeCharacter{03C9}{\ensuremath{\omega}}      % ω
\DeclareUnicodeCharacter{2013}{--}              % – (en dash)
\DeclareUnicodeCharacter{2014}{---}             % — (em dash)
\DeclareUnicodeCharacter{2018}{`}               % ‘
\DeclareUnicodeCharacter{2019}{'}               % ’
\DeclareUnicodeCharacter{201C}{``}              % “
\DeclareUnicodeCharacter{201D}{''}              % ”
\DeclareUnicodeCharacter{2070}{\ensuremath{^{0}}}        % ⁰
\DeclareUnicodeCharacter{2071}{\ensuremath{^{i}}}        % ⁱ
\DeclareUnicodeCharacter{2122}{\textsuperscript{TM}}     % ™
\DeclareUnicodeCharacter{2126}{\ensuremath{\Omega}}      % Ω (ohm sign)
\DeclareUnicodeCharacter{210F}{\ensuremath{\hbar}}       % ℏ
\DeclareUnicodeCharacter{2115}{\ensuremath{\mathbb{N}}}  % ℕ
\DeclareUnicodeCharacter{2102}{\ensuremath{\mathbb{C}}}  % ℂ
\DeclareUnicodeCharacter{211D}{\ensuremath{\mathbb{R}}}  % ℝ
\DeclareUnicodeCharacter{2124}{\ensuremath{\mathbb{Z}}}  % ℤ
\DeclareUnicodeCharacter{211A}{\ensuremath{\mathbb{Q}}}  % ℚ
\DeclareUnicodeCharacter{2192}{\ensuremath{\to}}         % →
\DeclareUnicodeCharacter{21D2}{\ensuremath{\Rightarrow}} % ⇒
\DeclareUnicodeCharacter{2200}{\ensuremath{\forall}}     % ∀
\DeclareUnicodeCharacter{2203}{\ensuremath{\exists}}     % ∃
\DeclareUnicodeCharacter{2208}{\ensuremath{\in}}         % ∈
\DeclareUnicodeCharacter{2209}{\ensuremath{\notin}}      % ∉
\DeclareUnicodeCharacter{2227}{\ensuremath{\wedge}}      % ∧
\DeclareUnicodeCharacter{2228}{\ensuremath{\vee}}        % ∨
\DeclareUnicodeCharacter{2229}{\ensuremath{\cap}}        % ∩
\DeclareUnicodeCharacter{222A}{\ensuremath{\cup}}        % ∪
\DeclareUnicodeCharacter{2245}{\ensuremath{\cong}}       % ≅
\DeclareUnicodeCharacter{2248}{\ensuremath{\approx}}     % ≈
\DeclareUnicodeCharacter{2260}{\ensuremath{\neq}}        % ≠
\DeclareUnicodeCharacter{2264}{\ensuremath{\leq}}        % ≤
\DeclareUnicodeCharacter{2265}{\ensuremath{\geq}}        % ≥
\DeclareUnicodeCharacter{22C5}{\ensuremath{\cdot}}       % ⋅
\DeclareUnicodeCharacter{2295}{\ensuremath{\oplus}}      % ⊕
\DeclareUnicodeCharacter{2297}{\ensuremath{\otimes}}     % ⊗

% ---- Mathematics & symbols ------------------------------------------
\usepackage{amsmath,amssymb,bm,mathrsfs,mathtools}
\usepackage{amsthm}             % theorem environments (load BEFORE hyperref)

% ---- Theorem environments (shared counter set, TECT canonical) -------
% All theorems / lemmas / propositions / corollaries / remarks share a
% single counter so cross-references read consistently across a paper.
% Definitions get a separate counter (definitions are numbered in their
% own sequence by editorial convention).
\newtheorem{theorem}{Theorem}
\newtheorem{lemma}[theorem]{Lemma}
\newtheorem{proposition}[theorem]{Proposition}
\newtheorem{corollary}[theorem]{Corollary}

\theoremstyle{remark}
\newtheorem{remark}[theorem]{Remark}
\newtheorem{example}[theorem]{Example}

\theoremstyle{definition}
\newtheorem{definition}[theorem]{Definition}
\newtheorem{conjecture}[theorem]{Conjecture}
\newtheorem{hypothesis}[theorem]{Hypothesis}

% ---- Tables / floats / graphics --------------------------------------
\usepackage{booktabs}           % \toprule / \midrule / \bottomrule
\usepackage{array}
\usepackage{graphicx}
\usepackage{enumitem}           % \begin{itemize}[leftmargin=...] options
\usepackage{hyphenat}           % \hyp{} for breakable hyphens in \texttt
\usepackage{seqsplit}           % \seqsplit{} for long unbreakable strings
\sloppy                         % allow stretchier inter-word spacing to
                                % reduce overfull \hbox warnings on long URLs
                                % and long \texttt tokens (paragraph-level).
\emergencystretch=3em           % last-resort hbox stretch budget

% ---- Cross-references (hyperref last among the standard packages,
%      cleveref AFTER hyperref) ----------------------------------------
\usepackage[%
  colorlinks=true,%
  linkcolor=blue,%
  citecolor=blue,%
  urlcolor=blue,%
  breaklinks=true,%
  bookmarks=true,%
  bookmarksnumbered=true%
]{hyperref}
\usepackage[capitalise,nameinlink]{cleveref}

% ---- TECT-canonical author block macros -----------------------------
% (defined here so every paper has the same author / affiliation style)
\newcommand{\tectauthor}{%
  \author{Jusang Lee}%
  \email{jtkor@outlook.com}%
  \affiliation{Independent researcher, Republic of Korea}%
}

% ---- TECT-canonical archive footer ----------------------------------
\newcommand{\tectarchivefooter}{%
  \par\medskip
  \noindent\small\textit{Canonical archive}:
  \url{https://github.com/TECT-OpenPhysics/TECT}.\par%
}

% ---- TECT TODO / AUDIT-FLAG / HONEST-SCOPE inline marks --------------
% Used in drafts to highlight pending audit items without disturbing
% the body text style. Suppress via \tectsuppressmarks (define empty).
\providecommand{\tectauditflag}[1]{\textcolor{red}{\textbf{[AUDIT-FLAG: #1]}}}
\providecommand{\tecttodo}[1]{\textcolor{orange}{\textbf{[TODO: #1]}}}
\providecommand{\tecthonestscope}[1]{\textit{Honest scope: #1.}}

% =====================================================================
% End of TECT canonical preamble
% =====================================================================

%   \begin{document}
%   ...
%   \end{document}
%
% Build (every paper): `pdflatex Paper-NN && pdflatex Paper-NN`
% (the second pass resolves \ref{} / \cref{} forward references).
% A bash/PowerShell wrapper that automates this is at
% Docs/papers/papers/_shared/build-paper.sh / build-paper.ps1.
% =====================================================================

% ---- Document class --------------------------------------------------
% PRD class option of revtex4-2 with [reprint,onecolumn] gives the same
% APS heading / by-line / abstract style as PRL but in a wider single
% column that handles long display equations and theorem environments
% without column-break artefacts. nofootinbib keeps citations out of
% footnotes. The 11pt option restores standard font metrics; with T1
% fontenc + lmodern this gives non-bitmap glyphs.
\documentclass[%
  aps,prd,reprint,onecolumn,nofootinbib,11pt,%
  amsmath,amssymb,floatfix,longbibliography%
]{revtex4-2}

% ---- Encoding & fonts ------------------------------------------------
\usepackage[T1]{fontenc}        % proper 8-bit font encoding (no font corruption)
\usepackage[utf8]{inputenc}     % allow UTF-8 source
\usepackage{lmodern}            % scalable Latin Modern (replaces bitmap CM)
\usepackage{textcomp}           % extra text symbols
\usepackage{microtype}          % subtle kerning improvements

% ---- UTF-8 → LaTeX character mappings --------------------------------
% Maps the Unicode characters that occur in TECT body text (legacy from
% the chat-content auto-archival rule, CLAUDE.md §4) to LaTeX-safe
% commands. Without these declarations, T1+inputenc rejects the chars
% with "Unicode character X not set up for use with LaTeX".
% Adding declarations centrally (here) is preferable to per-file edits.
\DeclareUnicodeCharacter{00A7}{\S}              % §  (section sign)
\DeclareUnicodeCharacter{00B2}{\ensuremath{^{2}}}        % ²
\DeclareUnicodeCharacter{00B3}{\ensuremath{^{3}}}        % ³
\DeclareUnicodeCharacter{00B5}{\ensuremath{\mu}}         % µ (micro)
\DeclareUnicodeCharacter{00B7}{\ensuremath{\cdot}}       % · (middle dot)
\DeclareUnicodeCharacter{00D7}{\ensuremath{\times}}      % ×
\DeclareUnicodeCharacter{00E4}{\"a}             % ä
\DeclareUnicodeCharacter{00E9}{\'e}             % é
\DeclareUnicodeCharacter{00F6}{\"o}             % ö
\DeclareUnicodeCharacter{00FC}{\"u}             % ü
\DeclareUnicodeCharacter{010C}{\v{C}}           % Č
\DeclareUnicodeCharacter{010D}{\v{c}}           % č
\DeclareUnicodeCharacter{0394}{\ensuremath{\Delta}}      % Δ
\DeclareUnicodeCharacter{03A3}{\ensuremath{\Sigma}}      % Σ
\DeclareUnicodeCharacter{03A9}{\ensuremath{\Omega}}      % Ω
\DeclareUnicodeCharacter{03B1}{\ensuremath{\alpha}}      % α
\DeclareUnicodeCharacter{03B2}{\ensuremath{\beta}}       % β
\DeclareUnicodeCharacter{03B3}{\ensuremath{\gamma}}      % γ
\DeclareUnicodeCharacter{03B4}{\ensuremath{\delta}}      % δ
\DeclareUnicodeCharacter{03B5}{\ensuremath{\varepsilon}} % ε
\DeclareUnicodeCharacter{03BA}{\ensuremath{\kappa}}      % κ
\DeclareUnicodeCharacter{03BB}{\ensuremath{\lambda}}     % λ
\DeclareUnicodeCharacter{03BC}{\ensuremath{\mu}}         % μ
\DeclareUnicodeCharacter{03BD}{\ensuremath{\nu}}         % ν
\DeclareUnicodeCharacter{03C0}{\ensuremath{\pi}}         % π
\DeclareUnicodeCharacter{03C1}{\ensuremath{\rho}}        % ρ
\DeclareUnicodeCharacter{03C3}{\ensuremath{\sigma}}      % σ
\DeclareUnicodeCharacter{03C4}{\ensuremath{\tau}}        % τ
\DeclareUnicodeCharacter{03C6}{\ensuremath{\varphi}}     % φ
\DeclareUnicodeCharacter{03C7}{\ensuremath{\chi}}        % χ
\DeclareUnicodeCharacter{03C8}{\ensuremath{\psi}}        % ψ
\DeclareUnicodeCharacter{03C9}{\ensuremath{\omega}}      % ω
\DeclareUnicodeCharacter{2013}{--}              % – (en dash)
\DeclareUnicodeCharacter{2014}{---}             % — (em dash)
\DeclareUnicodeCharacter{2018}{`}               % ‘
\DeclareUnicodeCharacter{2019}{'}               % ’
\DeclareUnicodeCharacter{201C}{``}              % “
\DeclareUnicodeCharacter{201D}{''}              % ”
\DeclareUnicodeCharacter{2070}{\ensuremath{^{0}}}        % ⁰
\DeclareUnicodeCharacter{2071}{\ensuremath{^{i}}}        % ⁱ
\DeclareUnicodeCharacter{2122}{\textsuperscript{TM}}     % ™
\DeclareUnicodeCharacter{2126}{\ensuremath{\Omega}}      % Ω (ohm sign)
\DeclareUnicodeCharacter{210F}{\ensuremath{\hbar}}       % ℏ
\DeclareUnicodeCharacter{2115}{\ensuremath{\mathbb{N}}}  % ℕ
\DeclareUnicodeCharacter{2102}{\ensuremath{\mathbb{C}}}  % ℂ
\DeclareUnicodeCharacter{211D}{\ensuremath{\mathbb{R}}}  % ℝ
\DeclareUnicodeCharacter{2124}{\ensuremath{\mathbb{Z}}}  % ℤ
\DeclareUnicodeCharacter{211A}{\ensuremath{\mathbb{Q}}}  % ℚ
\DeclareUnicodeCharacter{2192}{\ensuremath{\to}}         % →
\DeclareUnicodeCharacter{21D2}{\ensuremath{\Rightarrow}} % ⇒
\DeclareUnicodeCharacter{2200}{\ensuremath{\forall}}     % ∀
\DeclareUnicodeCharacter{2203}{\ensuremath{\exists}}     % ∃
\DeclareUnicodeCharacter{2208}{\ensuremath{\in}}         % ∈
\DeclareUnicodeCharacter{2209}{\ensuremath{\notin}}      % ∉
\DeclareUnicodeCharacter{2227}{\ensuremath{\wedge}}      % ∧
\DeclareUnicodeCharacter{2228}{\ensuremath{\vee}}        % ∨
\DeclareUnicodeCharacter{2229}{\ensuremath{\cap}}        % ∩
\DeclareUnicodeCharacter{222A}{\ensuremath{\cup}}        % ∪
\DeclareUnicodeCharacter{2245}{\ensuremath{\cong}}       % ≅
\DeclareUnicodeCharacter{2248}{\ensuremath{\approx}}     % ≈
\DeclareUnicodeCharacter{2260}{\ensuremath{\neq}}        % ≠
\DeclareUnicodeCharacter{2264}{\ensuremath{\leq}}        % ≤
\DeclareUnicodeCharacter{2265}{\ensuremath{\geq}}        % ≥
\DeclareUnicodeCharacter{22C5}{\ensuremath{\cdot}}       % ⋅
\DeclareUnicodeCharacter{2295}{\ensuremath{\oplus}}      % ⊕
\DeclareUnicodeCharacter{2297}{\ensuremath{\otimes}}     % ⊗

% ---- Mathematics & symbols ------------------------------------------
\usepackage{amsmath,amssymb,bm,mathrsfs,mathtools}
\usepackage{amsthm}             % theorem environments (load BEFORE hyperref)

% ---- Theorem environments (shared counter set, TECT canonical) -------
% All theorems / lemmas / propositions / corollaries / remarks share a
% single counter so cross-references read consistently across a paper.
% Definitions get a separate counter (definitions are numbered in their
% own sequence by editorial convention).
\newtheorem{theorem}{Theorem}
\newtheorem{lemma}[theorem]{Lemma}
\newtheorem{proposition}[theorem]{Proposition}
\newtheorem{corollary}[theorem]{Corollary}

\theoremstyle{remark}
\newtheorem{remark}[theorem]{Remark}
\newtheorem{example}[theorem]{Example}

\theoremstyle{definition}
\newtheorem{definition}[theorem]{Definition}
\newtheorem{conjecture}[theorem]{Conjecture}
\newtheorem{hypothesis}[theorem]{Hypothesis}

% ---- Tables / floats / graphics --------------------------------------
\usepackage{booktabs}           % \toprule / \midrule / \bottomrule
\usepackage{array}
\usepackage{graphicx}
\usepackage{enumitem}           % \begin{itemize}[leftmargin=...] options
\usepackage{hyphenat}           % \hyp{} for breakable hyphens in \texttt
\usepackage{seqsplit}           % \seqsplit{} for long unbreakable strings
\sloppy                         % allow stretchier inter-word spacing to
                                % reduce overfull \hbox warnings on long URLs
                                % and long \texttt tokens (paragraph-level).
\emergencystretch=3em           % last-resort hbox stretch budget

% ---- Cross-references (hyperref last among the standard packages,
%      cleveref AFTER hyperref) ----------------------------------------
\usepackage[%
  colorlinks=true,%
  linkcolor=blue,%
  citecolor=blue,%
  urlcolor=blue,%
  breaklinks=true,%
  bookmarks=true,%
  bookmarksnumbered=true%
]{hyperref}
\usepackage[capitalise,nameinlink]{cleveref}

% ---- TECT-canonical author block macros -----------------------------
% (defined here so every paper has the same author / affiliation style)
\newcommand{\tectauthor}{%
  \author{Jusang Lee}%
  \email{jtkor@outlook.com}%
  \affiliation{Independent researcher, Republic of Korea}%
}

% ---- TECT-canonical archive footer ----------------------------------
\newcommand{\tectarchivefooter}{%
  \par\medskip
  \noindent\small\textit{Canonical archive}:
  \url{https://github.com/TECT-OpenPhysics/TECT}.\par%
}

% ---- TECT TODO / AUDIT-FLAG / HONEST-SCOPE inline marks --------------
% Used in drafts to highlight pending audit items without disturbing
% the body text style. Suppress via \tectsuppressmarks (define empty).
\providecommand{\tectauditflag}[1]{\textcolor{red}{\textbf{[AUDIT-FLAG: #1]}}}
\providecommand{\tecttodo}[1]{\textcolor{orange}{\textbf{[TODO: #1]}}}
\providecommand{\tecthonestscope}[1]{\textit{Honest scope: #1.}}

% =====================================================================
% End of TECT canonical preamble
% =====================================================================

%   \begin{document}
%   ...
%   \end{document}
%
% Build (every paper): `pdflatex Paper-NN && pdflatex Paper-NN`
% (the second pass resolves \ref{} / \cref{} forward references).
% A bash/PowerShell wrapper that automates this is at
% Docs/papers/papers/_shared/build-paper.sh / build-paper.ps1.
% =====================================================================

% ---- Document class --------------------------------------------------
% PRD class option of revtex4-2 with [reprint,onecolumn] gives the same
% APS heading / by-line / abstract style as PRL but in a wider single
% column that handles long display equations and theorem environments
% without column-break artefacts. nofootinbib keeps citations out of
% footnotes. The 11pt option restores standard font metrics; with T1
% fontenc + lmodern this gives non-bitmap glyphs.
\documentclass[%
  aps,prd,reprint,onecolumn,nofootinbib,11pt,%
  amsmath,amssymb,floatfix,longbibliography%
]{revtex4-2}

% ---- Encoding & fonts ------------------------------------------------
\usepackage[T1]{fontenc}        % proper 8-bit font encoding (no font corruption)
\usepackage[utf8]{inputenc}     % allow UTF-8 source
\usepackage{lmodern}            % scalable Latin Modern (replaces bitmap CM)
\usepackage{textcomp}           % extra text symbols
\usepackage{microtype}          % subtle kerning improvements

% ---- UTF-8 → LaTeX character mappings --------------------------------
% Maps the Unicode characters that occur in TECT body text (legacy from
% the chat-content auto-archival rule, CLAUDE.md §4) to LaTeX-safe
% commands. Without these declarations, T1+inputenc rejects the chars
% with "Unicode character X not set up for use with LaTeX".
% Adding declarations centrally (here) is preferable to per-file edits.
\DeclareUnicodeCharacter{00A7}{\S}              % §  (section sign)
\DeclareUnicodeCharacter{00B2}{\ensuremath{^{2}}}        % ²
\DeclareUnicodeCharacter{00B3}{\ensuremath{^{3}}}        % ³
\DeclareUnicodeCharacter{00B5}{\ensuremath{\mu}}         % µ (micro)
\DeclareUnicodeCharacter{00B7}{\ensuremath{\cdot}}       % · (middle dot)
\DeclareUnicodeCharacter{00D7}{\ensuremath{\times}}      % ×
\DeclareUnicodeCharacter{00E4}{\"a}             % ä
\DeclareUnicodeCharacter{00E9}{\'e}             % é
\DeclareUnicodeCharacter{00F6}{\"o}             % ö
\DeclareUnicodeCharacter{00FC}{\"u}             % ü
\DeclareUnicodeCharacter{010C}{\v{C}}           % Č
\DeclareUnicodeCharacter{010D}{\v{c}}           % č
\DeclareUnicodeCharacter{0394}{\ensuremath{\Delta}}      % Δ
\DeclareUnicodeCharacter{03A3}{\ensuremath{\Sigma}}      % Σ
\DeclareUnicodeCharacter{03A9}{\ensuremath{\Omega}}      % Ω
\DeclareUnicodeCharacter{03B1}{\ensuremath{\alpha}}      % α
\DeclareUnicodeCharacter{03B2}{\ensuremath{\beta}}       % β
\DeclareUnicodeCharacter{03B3}{\ensuremath{\gamma}}      % γ
\DeclareUnicodeCharacter{03B4}{\ensuremath{\delta}}      % δ
\DeclareUnicodeCharacter{03B5}{\ensuremath{\varepsilon}} % ε
\DeclareUnicodeCharacter{03BA}{\ensuremath{\kappa}}      % κ
\DeclareUnicodeCharacter{03BB}{\ensuremath{\lambda}}     % λ
\DeclareUnicodeCharacter{03BC}{\ensuremath{\mu}}         % μ
\DeclareUnicodeCharacter{03BD}{\ensuremath{\nu}}         % ν
\DeclareUnicodeCharacter{03C0}{\ensuremath{\pi}}         % π
\DeclareUnicodeCharacter{03C1}{\ensuremath{\rho}}        % ρ
\DeclareUnicodeCharacter{03C3}{\ensuremath{\sigma}}      % σ
\DeclareUnicodeCharacter{03C4}{\ensuremath{\tau}}        % τ
\DeclareUnicodeCharacter{03C6}{\ensuremath{\varphi}}     % φ
\DeclareUnicodeCharacter{03C7}{\ensuremath{\chi}}        % χ
\DeclareUnicodeCharacter{03C8}{\ensuremath{\psi}}        % ψ
\DeclareUnicodeCharacter{03C9}{\ensuremath{\omega}}      % ω
\DeclareUnicodeCharacter{2013}{--}              % – (en dash)
\DeclareUnicodeCharacter{2014}{---}             % — (em dash)
\DeclareUnicodeCharacter{2018}{`}               % ‘
\DeclareUnicodeCharacter{2019}{'}               % ’
\DeclareUnicodeCharacter{201C}{``}              % “
\DeclareUnicodeCharacter{201D}{''}              % ”
\DeclareUnicodeCharacter{2070}{\ensuremath{^{0}}}        % ⁰
\DeclareUnicodeCharacter{2071}{\ensuremath{^{i}}}        % ⁱ
\DeclareUnicodeCharacter{2122}{\textsuperscript{TM}}     % ™
\DeclareUnicodeCharacter{2126}{\ensuremath{\Omega}}      % Ω (ohm sign)
\DeclareUnicodeCharacter{210F}{\ensuremath{\hbar}}       % ℏ
\DeclareUnicodeCharacter{2115}{\ensuremath{\mathbb{N}}}  % ℕ
\DeclareUnicodeCharacter{2102}{\ensuremath{\mathbb{C}}}  % ℂ
\DeclareUnicodeCharacter{211D}{\ensuremath{\mathbb{R}}}  % ℝ
\DeclareUnicodeCharacter{2124}{\ensuremath{\mathbb{Z}}}  % ℤ
\DeclareUnicodeCharacter{211A}{\ensuremath{\mathbb{Q}}}  % ℚ
\DeclareUnicodeCharacter{2192}{\ensuremath{\to}}         % →
\DeclareUnicodeCharacter{21D2}{\ensuremath{\Rightarrow}} % ⇒
\DeclareUnicodeCharacter{2200}{\ensuremath{\forall}}     % ∀
\DeclareUnicodeCharacter{2203}{\ensuremath{\exists}}     % ∃
\DeclareUnicodeCharacter{2208}{\ensuremath{\in}}         % ∈
\DeclareUnicodeCharacter{2209}{\ensuremath{\notin}}      % ∉
\DeclareUnicodeCharacter{2227}{\ensuremath{\wedge}}      % ∧
\DeclareUnicodeCharacter{2228}{\ensuremath{\vee}}        % ∨
\DeclareUnicodeCharacter{2229}{\ensuremath{\cap}}        % ∩
\DeclareUnicodeCharacter{222A}{\ensuremath{\cup}}        % ∪
\DeclareUnicodeCharacter{2245}{\ensuremath{\cong}}       % ≅
\DeclareUnicodeCharacter{2248}{\ensuremath{\approx}}     % ≈
\DeclareUnicodeCharacter{2260}{\ensuremath{\neq}}        % ≠
\DeclareUnicodeCharacter{2264}{\ensuremath{\leq}}        % ≤
\DeclareUnicodeCharacter{2265}{\ensuremath{\geq}}        % ≥
\DeclareUnicodeCharacter{22C5}{\ensuremath{\cdot}}       % ⋅
\DeclareUnicodeCharacter{2295}{\ensuremath{\oplus}}      % ⊕
\DeclareUnicodeCharacter{2297}{\ensuremath{\otimes}}     % ⊗

% ---- Mathematics & symbols ------------------------------------------
\usepackage{amsmath,amssymb,bm,mathrsfs,mathtools}
\usepackage{amsthm}             % theorem environments (load BEFORE hyperref)

% ---- Theorem environments (shared counter set, TECT canonical) -------
% All theorems / lemmas / propositions / corollaries / remarks share a
% single counter so cross-references read consistently across a paper.
% Definitions get a separate counter (definitions are numbered in their
% own sequence by editorial convention).
\newtheorem{theorem}{Theorem}
\newtheorem{lemma}[theorem]{Lemma}
\newtheorem{proposition}[theorem]{Proposition}
\newtheorem{corollary}[theorem]{Corollary}

\theoremstyle{remark}
\newtheorem{remark}[theorem]{Remark}
\newtheorem{example}[theorem]{Example}

\theoremstyle{definition}
\newtheorem{definition}[theorem]{Definition}
\newtheorem{conjecture}[theorem]{Conjecture}
\newtheorem{hypothesis}[theorem]{Hypothesis}

% ---- Tables / floats / graphics --------------------------------------
\usepackage{booktabs}           % \toprule / \midrule / \bottomrule
\usepackage{array}
\usepackage{graphicx}
\usepackage{enumitem}           % \begin{itemize}[leftmargin=...] options
\usepackage{hyphenat}           % \hyp{} for breakable hyphens in \texttt
\usepackage{seqsplit}           % \seqsplit{} for long unbreakable strings
\sloppy                         % allow stretchier inter-word spacing to
                                % reduce overfull \hbox warnings on long URLs
                                % and long \texttt tokens (paragraph-level).
\emergencystretch=3em           % last-resort hbox stretch budget

% ---- Cross-references (hyperref last among the standard packages,
%      cleveref AFTER hyperref) ----------------------------------------
\usepackage[%
  colorlinks=true,%
  linkcolor=blue,%
  citecolor=blue,%
  urlcolor=blue,%
  breaklinks=true,%
  bookmarks=true,%
  bookmarksnumbered=true%
]{hyperref}
\usepackage[capitalise,nameinlink]{cleveref}

% ---- TECT-canonical author block macros -----------------------------
% (defined here so every paper has the same author / affiliation style)
\newcommand{\tectauthor}{%
  \author{Jusang Lee}%
  \email{jtkor@outlook.com}%
  \affiliation{Independent researcher, Republic of Korea}%
}

% ---- TECT-canonical archive footer ----------------------------------
\newcommand{\tectarchivefooter}{%
  \par\medskip
  \noindent\small\textit{Canonical archive}:
  \url{https://github.com/TECT-OpenPhysics/TECT}.\par%
}

% ---- TECT TODO / AUDIT-FLAG / HONEST-SCOPE inline marks --------------
% Used in drafts to highlight pending audit items without disturbing
% the body text style. Suppress via \tectsuppressmarks (define empty).
\providecommand{\tectauditflag}[1]{\textcolor{red}{\textbf{[AUDIT-FLAG: #1]}}}
\providecommand{\tecttodo}[1]{\textcolor{orange}{\textbf{[TODO: #1]}}}
\providecommand{\tecthonestscope}[1]{\textit{Honest scope: #1.}}

% =====================================================================
% End of TECT canonical preamble
% =====================================================================
